\documentclass[12pt, twoside]{article}
\input{../preamble}

\fancyhead[LO]{La Scuola d'Italia / Dr.\ Huson / IB Math \\* 29 May 2026}
\fancyhead[RO]{\fbox{\begin{minipage}{6cm}\footnotesize First \& Last name:\\[0.5cm]\end{minipage}}}
\fancyfoot[C]{\thepage}

\begin{document}
\subsubsection*{8.6 Early Finishers: Transformations of Quadratics and a Cubic}

\noindent\emph{Answer on separate lined or graph paper, showing full working.
A graphing calculator is allowed to check your sketches, but show the algebra.
The parent functions are $f(x)=x^2$ and $f(x)=x^3$.}

\medskip

\begin{enumerate}[itemsep=0.55cm]

%--- Problem 1: quadratic in vertex form — describe, analyze, expand, sketch (freshly authored) ---
\item[1.] Let $f(x)=x^2$. The function $g$ is defined by
$\displaystyle g(x)=-2(x+3)^2+5$.

\begin{enumerate}[label=(\alph*), itemsep=0.3cm]
  \item Describe fully the sequence of transformations that maps the graph of
  $f$ onto the graph of $g$.
  \item Write down the coordinates of the vertex of $g$ and the equation of its
  axis of symmetry.
  \item State the range of $g$.
  \item Express $g(x)$ in the form $ax^2+bx+c$.
  \item Sketch the graph of $g$, labelling the vertex and the $y$-intercept.
\end{enumerate}

%--- Problem 2: build a quadratic from a description; commutativity of vertical scalings ---
\item[2.] The graph of $y=x^2$ undergoes the following transformations, in order:
a vertical stretch with scale factor $\tfrac{1}{2}$, then a reflection in the
$x$-axis, then a translation by the vector
$\left(\begin{smallmatrix}4\\-3\end{smallmatrix}\right)$.

\begin{enumerate}[label=(\alph*), itemsep=0.3cm]
  \item Find an equation for the resulting function $h$ in vertex form.
  \item Write down the coordinates of the vertex of $y=h(x)$, and determine
  whether the graph has any $x$-intercepts, justifying your answer.
  \item Does the order of the stretch and the reflection matter? Justify your
  answer.
\end{enumerate}

\newpage
%--- Problem 3: cubic with stretch + translation; point of inflection, intercepts, sketch ---
\item[3.] Let $f(x)=x^3$. The function $g$ is defined by
$\displaystyle g(x)=2(x+1)^3-16$.

\begin{enumerate}[label=(\alph*), itemsep=0.3cm]
  \item Describe fully the sequence of transformations that maps the graph of
  $f$ onto the graph of $g$.
  \item Write down the coordinates of the point of inflection of the graph of $g$.
  \item Find the coordinates of the $x$-intercept and of the $y$-intercept of the
  graph of $g$.
  \item Sketch the graph of $g$, showing the point of inflection and both
  intercepts.
\end{enumerate}

%--- Problem 4: challenge — reverse a quadratic by completing the square; build a cubic ---
\item[4.] \textbf{Challenge.}

\begin{enumerate}[label=(\alph*), itemsep=0.3cm]
  \item The graph of $y=x^2$ is transformed into the graph of
  $\displaystyle p(x)=-x^2+8x-13$. By completing the square, describe fully the
  sequence of transformations, and state the coordinates of the vertex of $p$.
  \item A curve has equation $y=a(x-h)^3+k$. Its point of inflection is at
  $(-1,2)$, and it passes through the point $(1,18)$. Find the values of $a$,
  $h$, and $k$, and hence write down the equation of the curve.
\end{enumerate}

\end{enumerate}
\end{document}
